\begin{pastexam}[label={px:ntu-112-bime-ds-05}]{臺大生機所 112 學年度（A）第 5 題}
（問答）

(10\%) When implementing the data structure of \textbf{queue} with a \textbf{linked list}, which
of the following design is better, (a) or (b), when efficiency is considered? (a) insert at the
front of the linked list and remove from the end of the linked list; (b) insert at the end of
the linked list and remove from the front of the linked list; Please justify your answer.
\end{pastexam}

\begin{pastsoln}
答案：(b)。 以一般的單向鏈結串列，記住頭指標 $\mathit{head}$ 和尾指標 $\mathit{tail}$ 來比較
（資結書虛擬碼~6.5 就是 (b) 的做法）。 佇列的前端在哪一頭，取出就在哪一頭做。

\begin{center}
\begin{tabular}{lll}
\toprule
 & 在串列前端做 & 在串列尾端做 \\
\midrule
插入 & $O(1)$：新節點指向舊的 $\mathit{head}$ & $O(1)$：$\mathit{tail}.\mathit{next} \gets$ 新節點 \\
刪除 & $O(1)$：$\mathit{head} \gets \mathit{head}.\mathit{next}$ & $O(n)$：要找到倒數第二個節點 \\
\bottomrule
\end{tabular}
\end{center}

\begin{itemize}
\item (a) 在前端插入、從尾端刪除。 插入 $O(1)$；但刪除尾端節點後，要把\textbf{倒數第二個}節點的
$\mathit{next}$ 設成 \NIL，並把 $\mathit{tail}$ 改指向它。 單向串列只能往後走，找倒數第二個節點
要從 $\mathit{head}$ 一路走過去，$\Theta(n)$。 每次 dequeue 都 $\Theta(n)$。
\item (b) 在尾端插入、從前端刪除。 enqueue 用 $\mathit{tail}$ 直接接在後面，dequeue 把
$\mathit{head}$ 往後移一格，兩者都 $O(1)$，不需要找任何前一個節點。
\end{itemize}
所以 (b) 比較好：兩個操作都是常數時間，$m$ 次操作共 $O(m)$；(a) 最壞是 $\Theta(mn)$。

補充：(a) 改用雙向鏈結串列（每個節點多一個 $\mathit{prev}$ 指標）也能做到兩邊都 $O(1)$，但每個
節點多花一個指標的空間，每次操作也要多維護一個指標，仍然不如 (b) 用單向串列就辦到。
\end{pastsoln}
